\chapter{Values, Completeness and Blocks of the Spectral Hierarchy}
\label{chap:torres}

The bilayer operator, its moments, and the semigroup
\(\mathfrak S=\langle90,120\rangle\) are defined only once in the
\cref{sec:jerarquia-espectral-infinita-rev7}. This chapter publishes values,
proves reconstructive completeness, and separates the global blocks. The
same integers \(90,120\) also occur as cardinalities in the
hexad--octad incidence, but that coincidence does not transport the incidence
\(\iota_{6,8}\) to the angular chart or convert its configurations into
mass moments.

\section{First evaluations of the hierarchy}

The indices
\[
 90,120,180,210,240,270,360,420
\]
are a selection of initial heights of \(\mathfrak S\), since
\[
 180=2\cdot90,\quad210=90+120,\quad240=2\cdot120,
\]
\[
 270=3\cdot90,\quad360=3\cdot120,
 \quad420=2(90+120).
\]
The moments \(s_n=q_-^n-q_+^n\) and the generating interface
\(\mathcal K_{\mathrm{ang}}^{90,120}=(s_{90},s_{120})\) are used here with the
already fixed convention; they are not redefined. All values are
determined by \(A\), \(\Cstar\), and the reader's quadratic recurrence.

\begin{table}[htbp]
\centering
\caption{Oriented moments of the global hierarchy.}
\label{tab:torres}
\begin{tabular}{@{}rr@{}}
\toprule
\(n\)&\(s_n\)\\
\midrule
90  & \(2.9516943928982796\times10^{-4}\)\\
120 & \(1.9683511250294992\times10^{-5}\)\\
180 & \(8.7345871869508335\times10^{-8}\)\\
210 & \(5.8181313057291163\times10^{-9}\)\\
240 & \(3.8754674969305336\times10^{-10}\)\\
270 & \(2.5814553607509555\times10^{-11}\)\\
360 & \(7.6293257871406076\times10^{-15}\)\\
420 & \(3.3850663628348497\times10^{-17}\)\\
\bottomrule
\end{tabular}
\end{table}

\begin{proposition}[Global Character of Moments]
The quantities \(s_n\) of \cref{tab:torres} are invariants of the fixed
angular chart. A species contributes integer coordinates that weight those
invariants, but does not modify \(q_\pm\) or introduce a tower of its own.
\end{proposition}
\begin{proof}
Each \(s_n\) is a function of \(q_-\) and \(q_+\). The theorem of inversion of
channels proves that the joint map
\((A,\Cstar)\mapsto(q_+,q_-)\) is injective. No species label appears
in these definitions.
\end{proof}

\begin{theorem}[Reconstructive completeness of the hierarchy]
\label{thm:completitud-reconstructiva-torres}
Let \(n_1<n_2<\cdots\) be the increasing enumeration without repetitions of
\(\mathfrak S\). For
\[
 a_k:=s_{n_k}=q_-^{n_k}-q_+^{n_k}>0,
\]
every positive ratio \(\rho\in\RR_{>0}\) admits an absolutely
convergent expansion in the hierarchy:
\[
 \boxed{\rho=
 \exp\!\left(\sum_{k\geq1}\nu_k a_k\right)}
 \qquad\text{for some sequence }(\nu_k)_{k\geq1}\subset\ZZ.
\]
\end{theorem}

\begin{proof}
As \(0<q_+<q_-<1\), it holds
\[
 0<a_k<q_-^{n_k}.
\]
The \(n_k\) are multiplos of \(30\) no menores that \(90\); therefore,
\[
 \sum_{k\geq1}a_k
 \leq\sum_{m\geq3}q_-^{30m}
 =\frac{q_-^{90}}{1-q_-^{30}}<\infty,
 \qquad a_k\longrightarrow0.
\]
Fix \(r_0=\log\rho\) and a tie-breaking rule for the nearest integer.
Define recursively
\[
 \nu_k=\operatorname{nint}\!\left(\frac{r_{k-1}}{a_k}\right),
 \qquad
 r_k=r_{k-1}-\nu_k a_k.
\]
The property of the nearest integer gives
\[
 |r_k|\leq\frac{a_k}{2}\longrightarrow0.
\]
The telescoping identity
\[
 r_0=\sum_{j=1}^{N}\nu_j a_j+r_N
\]
implies \(r_0=\sum_{j\geq1}\nu_j a_j\). The convergence is absolute, since
\[
 \begin{aligned}
 \sum_{k\geq1}|\nu_k|a_k
 &=\sum_{k\geq1}|r_{k-1}-r_k|\\
 &\leq |r_0|+\frac{a_1}{2}
 +\frac12\sum_{k\geq2}(a_{k-1}+a_k)<\infty.
 \end{aligned}
\]
By exponentiating, the statement is obtained.
\end{proof}

\begin{remark}[Scope and negative control]
The theorem is a reconstructive surjectivity of the tower language onto \(\RR_{>0}\): it proves resolution of the codomain, not physical selection. The algorithm receives \(r_0=\log\rho\); if \(\rho\) is a target magnitude, its coefficients are therefore a target-dependent reconstruction. Prospective provenance requires another arrow,
\[
 \text{species}\dashrightarrow\text{extension}
 \longrightarrow\text{route}\longrightarrow\text{record}
 \longrightarrow\text{signature},
\]
and does not follow from this expansion. The conclusion would fail if the \(a_k\) did not
tend to zero, if the bound on the residue ceased to hold, or if the error series
were not summable.
\end{remark}

\section{Global blocks of mass action}

The mass action will use the triple
\[
 \Dfour,\qquad s_{120},\qquad \zCorr=135\Dfour^2.
\]
Their values are calculated only once. The coordinates
\(k,\nu_{120},\nu_{270}\) of a signature specify how many times
each block enters the exponent; they do not change its value.

It is useful to distinguish two objects that share the index \(270\):
\[
 \boxed{\zCorr=135\Dfour^2
 =1.66922699663643\times10^{-6}}
\]
and
\[
 \boxed{\sSpec=q_-^{270}-q_+^{270}
 =2.58145536075096\times10^{-11}.}
\]
The first is a quadratic invariant constructed from \(\Dfour\); the
second is a spectral moment of the semigroup. The coincidence of the subscript does not
authorize their identification.

\section{Higher-order combination}

The global combination
\[
 \delta\Dfour^{\rm tower}
 =\frac{s_{180}}{24}-\frac{s_{240}}{54}
 -24s_{360}+\frac76s_{420}
\]
evaluates a
\[
 \delta\Dfour^{\rm tower}
 =3.632051471908759123889070967\times10^{-9}.
\]
The comparison coordinate used in the selection of the
coefficients is defined as the scalar data
\[
 \delta\Dfour
 :=3.632051471908972784661641720\times10^{-9}.
\]
Consequently, the residue of the combination with respect to that coordinate is
\[
 \delta\Dfour^{\rm tower}-\delta\Dfour
 =-2.1366077257075\times10^{-22}.
\]
The small quantity is therefore a difference between two objects and not the value
of the combination. All the summands belong to the same hierarchy, and their
coefficients are global. Because the comparison coordinate intervened in the
selection of the coefficients, this equality is reconstructive. The
combination is not added again to the electronic formula that already uses
\(\Dfour\), since that substitution would duplicate the same correction.
